\section{Banked Cloze}
\subsection{Unit 1}

\begin{question}{Banked Cloze}

Through online interactions, young people can share ideas, acquire information, and be exposed to many opportunities out there. Social media gives them the avenue to make new friends and allows them to 
\textbf{\color{red} maintain} contact with their old friends.
 
\vspace{1em}

However, too much social media also leads to addiction. We are quite familiar with the \textbf{\color{red} scenario} of people sitting next to each other, not talking, but just staring at their phones. Young people are more likely to be \textbf{\color{red} immersed} in social media rather than spend time with their families and friends. They may check their social media accounts first thing in the morning and check social media updates last thing before going to sleep. With too many distractions, they cannot fully \textbf{\color{red} concentrate} on their tasks. 

\vspace{1em}

Social media also leads to \textbf{\color{red} obsession}. Some people are crazy about likes and comments and they will do anything to get them on a daily basis. Everything they do, they want to post on social media to get some form of \textbf{\color{red} approval} from strangers. They work tirelessly to buy the latest clothes, eat expensive meals, and visit unusual places, just for a show. They even do things that they would never do in real life just to impress people. Since they cannot get the same popularity in real life as they do in social media, they yield to the \textbf{\color{red} lure} of social networks. However, many of them live in \textbf{\color{red} continual} fear of being less popular. Everything falls apart when their popularity really \textbf{\color{red} diminishes}, and they may even get into serious psychological problems.

\vspace{1em}

The creators of social networks perhaps did not \textbf{\color{red} anticipate}
that their creations would bring these problems. Doing away with social media may not be practical, but some form of regulation should be put into place to ensure that people can really benefit from the technology.

\end{question}

\newpage

\subsection{Unit 2}

\begin{question}{Banked Cloze}

Zheng He was a Chinese explorer who led seven great voyages on behalf of the emperor of the Ming Dynasty. His voyages, which aimed to promote trade and \textbf{\color{red} diplomacy} with other states, lasted from 1405 to 1433.

\vspace{1em}

Zheng He was born around 1371 in Yunnan, southwestern China. He was talented in many \textbf{\color{red} realms}, especially in ship construction. In 1403, the emperor ordered the construction of the Treasure Fleet and \textbf{\color{red} appointed} Zheng He as the commander of this fleet. This began Zheng He's \textbf{\color{red} epic} voyages, some of the most impressive exploration journeys in history.

\vspace{1em}

Zheng He's first voyage began in July 1405. His fleet traveled to present-day Vietnam, Indonesia, Malaysia, and then crossed the Indian Ocean and sailed west to today's India. They traded with the local people for spices, ivory, \textbf{\color{red} exotic} animals, and other goods.

\vspace{1em}

During the next six expeditions, Zheng He made trade \textbf{\color{red} pacts} with the rulers of different states, brought back gifts, and established friendly relationships with them. He even brought back with him \textbf{\color{red} ambassadors} of various states to meet with the emperor of the Ming Dynasty.

\vspace{1em}

Zheng He's seven great voyages promoted the \textbf{\color{red} expansion} of China's cultural and economic influence in the world. He \textbf{\color{red} forged} new, friendly ties with many states on the Asian and African continents, while developing the trade relations with them. Zheng He's feat (壮举) was unprecedented in the history of ocean \textbf{\color{red} navigation}. He is still remembered today as one of the world's greatest navigators and is a source of pride for the whole Chinese nation.

\end{question}

\newpage

\subsection{Unit 3}

\begin{question}{Banked Cloze}

The benefits of travel go beyond making memories and meeting new people. Staying away from the \textbf{\color{red} trivial} things of daily life and exploring a new place can have a remarkably positive impact on your emotional well-being and your \textbf{\color{red} cognitive} powers.

\vspace{1em}

Travel can improve your emotional well-being in two aspects. First, changing your daily routine, for example, taking a trip to a nearby town on a weekend day, can help \textbf{\color{red} console} your mind and deal with the daily \textbf{\color{red} annoyances}. Second, physical exercise is known to improve emotional well-being and travel offers plenty of opportunities to get active. Getting to know a new destination by \textbf{\color{red} embracing} the great outdoors can boost energy levels and improve your mood.

\vspace{1em}

Travel can also improve your cognitive powers due to the exposure to different cultures. Whether you venture abroad or simply to the next town over, a trip away from home can help enhance your \textbf{\color{red} positivity} and reveal a world of possibilities beyond your daily perspective. Something as simple as learning a new \textbf{\color{red} recipe} or changing the way you spend your downtime can have a dramatic effect on your well-being. It's been proven that new experiences – particularly ones that allow you to \textbf{\color{red} immerse} yourself in a different culture – improve the performance of your brain, \textbf{\color{red} loosen} the restrictions, and increase creativity.

\vspace{1em}

Evidently, travel itself \textbf{\color{red} endows} you with the freedom to do what you love, to rest, and to try something new. For many people, travel is not simply an enjoyable recreation, but an essential part of maintaining a positive mentality and improving cognitive abilities.

\end{question}


\newpage

\subsection{Unit 4}

\begin{question}{Banked Cloze}

According to a study conducted in the U.K., four out of every five employees are happy at work. Surprisingly, the study revealed that the \textbf{\color{red} notion} that people are happy when they have a good salary, isn't true. So, what should people do to keep their spirits up and, at the same time, \textbf{\color{red} foster} a sense of joy at work? Here are two suggestions which are believed to be important in achieving happiness at work.

\vspace{1em}

Start with a positive outlook. Happiness is a state of mind, \textbf{\color{red} reflecting} an attitude. Staying happy at work is totally based on your motivation and a positive outlook toward your job, not on \textbf{\color{red} monetary} rewards or material gain. Always remembering the good aspects of the work and \textbf{\color{red} refraining} from thinking about what makes you unhappy are important to happiness. Negativity and \textbf{\color{red} complaining} about bad things are easy; it is looking at the bright side that is the difficult part of a job.

\vspace{1em}

Challenge yourself. Rather than looking for help and \textbf{\color{red} incentives} from other people, challenge yourself and take charge of your own growth professionally. Getting \textbf{\color{red} bored} with work is one of the primary factors that cause people to switch jobs. So find new challenges and deal with them by making \textbf{\color{red} exhaustive} plans and following them. It is a good idea to have help from your manager if you hit a \textbf{\color{red} barrier}, but completely relying on someone for your career advancement is unwise.

\vspace{1em}

Facing the problems and challenges at work and tackling them by yourself will give you a sense of accomplishment and help you achieve happiness at work.

\end{question}

\newpage

\subsection{Unit 5}

\begin{question}{Banked Cloze}

China's first attempt at independently developing its astronautic industry can be dated back to the 1950s, but it was not until the 1960s that any substantial progress was made.

\vspace{1em}

A major breakthrough in China's astronautic technology came on April 24, 1970, when China succeeded in launching its first satellite. This development marked a \textbf{\color{red} milestone} in China's research on man-made satellite.

\vspace{1em}

The 1980s was a period of \textbf{\color{red} soaring} growth for Chinese space technology. The successful launch of the Long March-3 brought the country into a leading position in carrier rocket technology around the world.

\vspace{1em}

In 1992, the Chinese government created a \textbf{\color{red} blueprint} and decided to implement its manned space program, which laid claim to being the largest and most technically complex \textbf{\color{red} endeavor} of China's space exploration in the late 20th century and early 21st century.

\vspace{1em}

In October 2003, the Shenzhou-5 manned spaceship \textbf{\color{red} ascended} into space and sent the first Chinese astronaut into orbit. This momentous space journey marked the realization of China's dream of a manned space flight.

\vspace{1em}

On October 16, 2021, Shenzhou-13 was launched and entered its \textbf{\color{red} designated} orbit successfully. After entering its orbit, the spaceship was 
\textbf{\color{red} docked} with the Tianhe core module. The crew completed a six-month stay, the longest one since the beginning of the program, and returned to the earth safely on April 16, 2022.

\vspace{1em}

China now is \textbf{\color{red} pushing} forward with its deep space exploration, and aims to contribute its wisdom to utilizing (利用) the outer space peacefully. The achievements of China's space exploration will
\textbf{\color{red} trigger} the interest of more young people who are willing to devote themselves to the 
\textbf{\color{red} lofty} cause. China's space exploration, of course, will start a new journey.

\end{question}

\newpage
\subsection{Unit 6}

\begin{question}{Banked Cloze}

While older people in China typically believe that saving is a virtue, their children might have other ideas. Many young people refuse to \textbf{\color{red} comply} with the traditional wisdom. They prefer to spend their money, buy now, and worry about paying later. According to a recent report, China's household savings rate has \textbf{\color{red} tumbled} in recent years.

\vspace{1em}

The younger generation is more willing to spend money on fine food, tourism, gymnasium and fitness \textbf{\color{red} sessions}, and high-tech products. Some of them spend all their salary by the end of each month, with part of the salary used to \textbf{\color{red} repay} loans.

\vspace{1em}

In recent years, Chinese banks and finance companies have launched various kinds of consumer credit services to encourage consumption. Unlike \textbf{\color{red} mortgages}, they don't require guarantees (抵押品). The services have quickly become popular with the younger generation, who are also the main force for mobile payment, and are willing to \textbf{\color{red} embrace} new technology and financial products. Generally, consumption on credit can be \textbf{\color{red} beneficial} to the development of the financial and retail sectors, and will eventually \textbf{\color{red} boost} economic growth.

\vspace{1em}

However, to ensure the healthy development of consumer finance, it is important for financial institutions to manage the default (违约) risks, especially for people who have \textbf{\color{red} burdensome} debts and may have difficulty in repaying their loans. It is also important to educate young people about the importance of maintaining a good credit record, especially for those easily \textbf{\color{red} tempted} by compulsive buying. They need to know a good credit record is important for their financial well-being.

\end{question}
